\label{sec:conclude}
AI agents increasingly rely on tool outputs as the substrate for reasoning, but existing MCP deployments frequently deliver large structured results through a JSON-RPC control plane that is ill-suited for tabular payloads. This paper presented \system, a dynamic format selection framework that splits MCP results into a small control-plane envelope (text plus a machine-readable descriptor) and a streamable HTTP data plane for large payload bytes. For tabular results, \system supports JSON records for compatibility and Parquet blob/stream delivery for efficient transfer, incremental consumption, and early termination; it further supports target-driven selection (min bytes, min latency, or min time-to-first-rows) and codec choices within Parquet.

Our evaluation on synthetic grids, TPC-DS slices, and BIRD dev query results shows that the design is practical and yields substantial reductions in aggregate wire volume versus JSON baselines, while preserving low millisecond-scale fetch latency on localhost. The results also highlight the importance of selection placement: client-driven selection benefits from explicit hints but incurs a ``describe'' overhead that can be removed via one-shot server-side selection.

\smallskip
\noindent\textbf{Limitations and future work.}
Our current benchmarks focus on localhost and fixed bandwidth assumptions; evaluating under realistic WAN conditions and varying congestion is an immediate next step. Second, while we treat unstructured outputs as first-class, richer multimodal payloads (e.g., images, embeddings, and mixed text+table outputs) require additional conventions for schema/typing and safe decoding. Third, selection policies can be made more data-driven: online learning can incorporate workload drift and privacy constraints, and codec selection can be extended beyond simple rule-based strategies. Finally, integrating end-to-end authorization, auditing, and privacy-preserving transformations into the data plane remains essential for multi-tenant streamable HTTP deployments.